% ============================================================
% Question Bank: ch05-02 Solving Two-Step Equations
% Topic Title: Solving Two-Step Equations
% CCSS: 7.EE.B.4.A
% Questions: 30 (18 MC + 7 SA + 5 Graphical)
% ============================================================

% --- Multiple Choice Questions (18) ---

\begin{practiceQuestion}{5-2-q01}{mc}
\begin{questionText}
What is the solution to $3x + 7 = 28$?
\end{questionText}
\choiceA{$x = 7$}
\choiceB{$x = 11\frac{2}{3}$}
\choiceC{$x = 10$}
\choiceD{$x = 35$}
\correctAnswer{A}
\explanation{Subtract $7$: $3x = 21$. Divide by $3$: $x = 7$. Check: $3(7) + 7 = 28$ \checkmark.}
\end{practiceQuestion}

\begin{practiceQuestion}{5-2-q02}{mc}
\begin{questionText}
Solve $5n - 8 = 27$.
\end{questionText}
\choiceA{$n = 3.8$}
\choiceB{$n = 7$}
\choiceC{$n = 19$}
\choiceD{$n = -7$}
\correctAnswer{B}
\explanation{Add $8$: $5n = 35$. Divide by $5$: $n = 7$. Check: $5(7) - 8 = 27$ \checkmark.}
\end{practiceQuestion}

\begin{practiceQuestion}{5-2-q03}{mc}
\begin{questionText}
What is the value of $y$ in $\dfrac{y}{4} + 6 = 11$?
\end{questionText}
\choiceA{$y = 68$}
\choiceB{$y = 44$}
\choiceC{$y = 20$}
\choiceD{$y = 1.25$}
\correctAnswer{C}
\explanation{Subtract $6$: $\frac{y}{4} = 5$. Multiply by $4$: $y = 20$. Check: $\frac{20}{4} + 6 = 5 + 6 = 11$ \checkmark.}
\end{practiceQuestion}

\begin{practiceQuestion}{5-2-q04}{mc}
\begin{questionText}
Solve $-2x + 9 = 1$.
\end{questionText}
\choiceA{$x = -4$}
\choiceB{$x = -5$}
\choiceC{$x = 4$}
\choiceD{$x = 5$}
\correctAnswer{C}
\explanation{Subtract $9$: $-2x = -8$. Divide by $-2$: $x = 4$. Check: $-2(4) + 9 = -8 + 9 = 1$ \checkmark.}
\end{practiceQuestion}

\begin{practiceQuestion}{5-2-q05}{mc}
\begin{questionText}
Solve $14 = 2m + 6$.
\end{questionText}
\choiceA{$m = 4$}
\choiceB{$m = 10$}
\choiceC{$m = 8$}
\choiceD{$m = 3$}
\correctAnswer{A}
\explanation{Subtract $6$: $8 = 2m$. Divide by $2$: $m = 4$. Check: $2(4) + 6 = 14$ \checkmark.}
\end{practiceQuestion}

\begin{practiceQuestion}{5-2-q06}{mc}
\begin{questionText}
What is the first step to solve $7x - 3 = 25$?
\end{questionText}
\choiceA{Divide both sides by $7$}
\choiceB{Subtract $3$ from both sides}
\choiceC{Add $3$ to both sides}
\choiceD{Multiply both sides by $7$}
\correctAnswer{C}
\explanation{To undo the subtraction of $3$, add $3$ to both sides first: $7x = 28$.}
\end{practiceQuestion}

\begin{practiceQuestion}{5-2-q07}{mc}
\begin{questionText}
Solve $\dfrac{n}{6} - 5 = 3$.
\end{questionText}
\choiceA{$n = 48$}
\choiceB{$n = -12$}
\choiceC{$n = 18$}
\choiceD{$n = 12$}
\correctAnswer{A}
\explanation{Add $5$: $\frac{n}{6} = 8$. Multiply by $6$: $n = 48$. Check: $\frac{48}{6} - 5 = 8 - 5 = 3$ \checkmark.}
\end{practiceQuestion}

\begin{practiceQuestion}{5-2-q08}{mc}
\begin{questionText}
A gym charges $\$12$ per month plus a $\$30$ registration fee. James paid $\$102$ total. How many months did he attend?
\end{questionText}
\choiceA{$6$ months}
\choiceB{$8$ months}
\choiceC{$11$ months}
\choiceD{$5$ months}
\correctAnswer{A}
\explanation{$12m + 30 = 102$. Subtract $30$: $12m = 72$. Divide by $12$: $m = 6$.}
\end{practiceQuestion}

\begin{practiceQuestion}{5-2-q09}{mc}
\begin{questionText}
Solve $-5k + 20 = -10$.
\end{questionText}
\choiceA{$k = 2$}
\choiceB{$k = -6$}
\choiceC{$k = 6$}
\choiceD{$k = -2$}
\correctAnswer{C}
\explanation{Subtract $20$: $-5k = -30$. Divide by $-5$: $k = 6$. Check: $-5(6) + 20 = -30 + 20 = -10$ \checkmark.}
\end{practiceQuestion}

\begin{practiceQuestion}{5-2-q10}{mc}
\begin{questionText}
Which equation has the solution $x = 5$?
\end{questionText}
\choiceA{$4x + 3 = 18$}
\choiceB{$4x + 3 = 23$}
\choiceC{$4x - 3 = 23$}
\choiceD{$4x - 5 = 25$}
\correctAnswer{B}
\explanation{Substitute $x = 5$: $4(5) + 3 = 20 + 3 = 23$. So $4x + 3 = 23$ is correct.}
\end{practiceQuestion}

\begin{practiceQuestion}{5-2-q11}{mc}
\begin{questionText}
Solve $9 + 3x = 30$.
\end{questionText}
\choiceA{$x = 13$}
\choiceB{$x = 3$}
\choiceC{$x = 7$}
\choiceD{$x = 10$}
\correctAnswer{C}
\explanation{Subtract $9$: $3x = 21$. Divide by $3$: $x = 7$. Check: $9 + 3(7) = 9 + 21 = 30$ \checkmark.}
\end{practiceQuestion}

\begin{practiceQuestion}{5-2-q12}{mc}
\begin{questionText}
Solve $\dfrac{w}{5} + 12 = 20$.
\end{questionText}
\choiceA{$w = 40$}
\choiceB{$w = 160$}
\choiceC{$w = 8$}
\choiceD{$w = 1.6$}
\correctAnswer{A}
\explanation{Subtract $12$: $\frac{w}{5} = 8$. Multiply by $5$: $w = 40$. Check: $\frac{40}{5} + 12 = 8 + 12 = 20$ \checkmark.}
\end{practiceQuestion}

\begin{practiceQuestion}{5-2-q13}{mc}
\begin{questionText}
You have $\$75$ in your bank account and withdraw $\$9$ each week. After how many weeks will you have $\$21$ left?
\end{questionText}
\choiceA{$5$ weeks}
\choiceB{$8$ weeks}
\choiceC{$7$ weeks}
\choiceD{$6$ weeks}
\correctAnswer{D}
\explanation{$75 - 9w = 21$. Subtract $75$: $-9w = -54$. Divide by $-9$: $w = 6$.}
\end{practiceQuestion}

\begin{practiceQuestion}{5-2-q14}{mc}
\begin{questionText}
Solve $-3x - 7 = 14$.
\end{questionText}
\choiceA{$x = -7$}
\choiceB{$x = 7$}
\choiceC{$x = -\dfrac{7}{3}$}
\choiceD{$x = \dfrac{7}{3}$}
\correctAnswer{A}
\explanation{Add $7$: $-3x = 21$. Divide by $-3$: $x = -7$. Check: $-3(-7) - 7 = 21 - 7 = 14$ \checkmark.}
\end{practiceQuestion}

\begin{practiceQuestion}{5-2-q15}{mc}
\begin{questionText}
Solve $8 = \dfrac{x}{3} + 2$.
\end{questionText}
\choiceA{$x = 30$}
\choiceB{$x = 18$}
\choiceC{$x = 2$}
\choiceD{$x = 6$}
\correctAnswer{B}
\explanation{Subtract $2$: $6 = \frac{x}{3}$. Multiply by $3$: $x = 18$. Check: $\frac{18}{3} + 2 = 6 + 2 = 8$ \checkmark.}
\end{practiceQuestion}

\begin{practiceQuestion}{5-2-q16}{mc}
\begin{questionText}
A student solved $4x + 10 = 26$ and got $x = 9$. What error did the student most likely make?
\end{questionText}
\choiceA{Added $10$ instead of subtracting}
\choiceB{Divided by $4$ before subtracting $10$}
\choiceC{Multiplied both sides by $4$}
\choiceD{Subtracted $10$ from the left side only}
\correctAnswer{A}
\explanation{If the student added $10$: $4x = 36$, then $x = 9$. The correct step is to subtract $10$: $4x = 16$, $x = 4$.}
\end{practiceQuestion}

\begin{practiceQuestion}{5-2-q17}{mc}
\begin{questionText}
Solve $6p + 5 = -7$.
\end{questionText}
\choiceA{$p = 2$}
\choiceB{$p = -\dfrac{1}{3}$}
\choiceC{$p = -2$}
\choiceD{$p = \dfrac{1}{3}$}
\correctAnswer{C}
\explanation{Subtract $5$: $6p = -12$. Divide by $6$: $p = -2$. Check: $6(-2) + 5 = -12 + 5 = -7$ \checkmark.}
\end{practiceQuestion}

\begin{practiceQuestion}{5-2-q18}{mc}
\begin{questionText}
Solve $50 - 4x = 18$.
\end{questionText}
\choiceA{$x = -8$}
\choiceB{$x = 17$}
\choiceC{$x = 8$}
\choiceD{$x = 12$}
\correctAnswer{C}
\explanation{Subtract $50$: $-4x = -32$. Divide by $-4$: $x = 8$. Check: $50 - 4(8) = 50 - 32 = 18$ \checkmark.}
\end{practiceQuestion}

% --- Short Answer Questions (7) ---

\begin{practiceQuestion}{5-2-q19}{sa}
\begin{questionText}
Solve $2x + 15 = 31$.
\end{questionText}
\correctAnswer{$x = 8$}
\explanation{Subtract $15$: $2x = 16$. Divide by $2$: $x = 8$. Check: $2(8) + 15 = 31$ \checkmark.}
\end{practiceQuestion}

\begin{practiceQuestion}{5-2-q20}{sa}
\begin{questionText}
Solve $\dfrac{y}{7} - 4 = 2$.
\end{questionText}
\correctAnswer{$y = 42$}
\explanation{Add $4$: $\frac{y}{7} = 6$. Multiply by $7$: $y = 42$. Check: $\frac{42}{7} - 4 = 6 - 4 = 2$ \checkmark.}
\end{practiceQuestion}

\begin{practiceQuestion}{5-2-q21}{sa}
\begin{questionText}
Solve $-8n + 3 = -37$.
\end{questionText}
\correctAnswer{$n = 5$}
\explanation{Subtract $3$: $-8n = -40$. Divide by $-8$: $n = 5$. Check: $-8(5) + 3 = -40 + 3 = -37$ \checkmark.}
\end{practiceQuestion}

\begin{practiceQuestion}{5-2-q22}{sa}
\begin{questionText}
A ride-share app charges $\$2.50$ per mile plus a $\$3.00$ base fee. The total fare was $\$15.50$. How many miles was the ride?
\end{questionText}
\correctAnswer{$5$ miles}
\explanation{$2.50m + 3 = 15.50$. Subtract $3$: $2.50m = 12.50$. Divide by $2.50$: $m = 5$.}
\end{practiceQuestion}

\begin{practiceQuestion}{5-2-q23}{sa}
\begin{questionText}
Solve $100 - 7x = 37$.
\end{questionText}
\correctAnswer{$x = 9$}
\explanation{Subtract $100$: $-7x = -63$. Divide by $-7$: $x = 9$. Check: $100 - 7(9) = 100 - 63 = 37$ \checkmark.}
\end{practiceQuestion}

\begin{practiceQuestion}{5-2-q24}{sa}
\begin{questionText}
Solve $\dfrac{k}{8} + 9 = 15$.
\end{questionText}
\correctAnswer{$k = 48$}
\explanation{Subtract $9$: $\frac{k}{8} = 6$. Multiply by $8$: $k = 48$. Check: $\frac{48}{8} + 9 = 6 + 9 = 15$ \checkmark.}
\end{practiceQuestion}

\begin{practiceQuestion}{5-2-q25}{sa}
\begin{questionText}
Solve $11 = 4x - 5$.
\end{questionText}
\correctAnswer{$x = 4$}
\explanation{Add $5$: $16 = 4x$. Divide by $4$: $x = 4$. Check: $4(4) - 5 = 16 - 5 = 11$ \checkmark.}
\end{practiceQuestion}

% --- Graphical Questions (3 GMC + 2 GSA) ---

\begin{practiceQuestion}{5-2-q26}{gmc}
\begin{questionText}
Look at the balance scale. Both sides are equal. What is the value of $x$?

\begin{center}
\begin{tikzpicture}[scale=0.7]
  % Fulcrum
  \draw[thick] (0,-0.5) -- (-0.4,-1.2) -- (0.4,-1.2) -- cycle;
  % Beam
  \draw[thick] (-4.5,0) -- (4.5,0);
  % Left pan
  \draw[thick] (-4,0) -- (-4.8,-0.6) -- (-3.2,-0.6) -- (-4,0);
  \draw[thick,fill=funBlue!20] (-4.8,-0.6) -- (-4.8,-1) -- (-3.2,-1) -- (-3.2,-0.6) -- cycle;
  % Left contents: 2 boxes and 4 units
  \node[draw,fill=funBlue!40,minimum size=0.4cm] at (-4.4,-1.4) {$x$};
  \node[draw,fill=funBlue!40,minimum size=0.4cm] at (-3.6,-1.4) {$x$};
  \node at (-4,-1.9) {$+ \; 4$};
  % Right pan
  \draw[thick] (4,0) -- (4.8,-0.6) -- (3.2,-0.6) -- (4,0);
  \draw[thick,fill=funOrange!20] (3.2,-0.6) -- (3.2,-1) -- (4.8,-1) -- (4.8,-0.6) -- cycle;
  \node at (4,-1.4) {$18$};
\end{tikzpicture}
\end{center}
\end{questionText}
\choiceA{$x = 9$}
\choiceB{$x = 7$}
\choiceC{$x = 14$}
\choiceD{$x = 11$}
\correctAnswer{B}
\explanation{The equation is $2x + 4 = 18$. Subtract $4$: $2x = 14$. Divide by $2$: $x = 7$.}
\end{practiceQuestion}

\begin{practiceQuestion}{5-2-q27}{gmc}
\begin{questionText}
The table shows a pattern. What is the value of $x$ when the output is $41$?

\begin{center}
\begin{tabular}{|c|c|}
\hline
\textbf{Input ($x$)} & \textbf{Output} \\
\hline
$1$ & $8$ \\
\hline
$2$ & $11$ \\
\hline
$3$ & $14$ \\
\hline
$?$ & $41$ \\
\hline
\end{tabular}
\end{center}
\end{questionText}
\choiceA{$x = 12$}
\choiceB{$x = 11$}
\choiceC{$x = 13$}
\choiceD{$x = 14$}
\correctAnswer{A}
\explanation{The pattern is $3x + 5$. Solve $3x + 5 = 41$: subtract $5$ to get $3x = 36$, divide by $3$: $x = 12$. Check: $3(12) + 5 = 41$ \checkmark.}
\end{practiceQuestion}

\begin{practiceQuestion}{5-2-q28}{gmc}
\begin{questionText}
The diagram shows a triangle with a perimeter of $39$ cm. What is the value of $x$?

\begin{center}
\begin{tikzpicture}[scale=0.8]
  \draw[thick,fill=funBlue!10] (0,0) -- (5,0) -- (2.5,3.5) -- cycle;
  \node[below] at (2.5,0) {$2x + 1$};
  \node[left] at (1,1.9) {$x + 3$};
  \node[right] at (4,1.9) {$x + 5$};
\end{tikzpicture}
\end{center}
\end{questionText}
\choiceA{$x = 7.5$}
\choiceB{$x = 10$}
\choiceC{$x = 6$}
\choiceD{$x = 8$}
\correctAnswer{A}
\explanation{Perimeter: $(2x + 1) + (x + 3) + (x + 5) = 39$. Combine: $4x + 9 = 39$. Subtract $9$: $4x = 30$. Divide by $4$: $x = 7.5$.}
\end{practiceQuestion}

\begin{practiceQuestion}{5-2-q29}{gsa}
\begin{questionText}
The bar model shows an equation. Solve for $x$.

\begin{center}
\begin{tikzpicture}[scale=0.7]
  % Top bar: total
  \draw[thick,fill=funOrange!20] (0,1.5) rectangle (10,2.5);
  \node at (5,2) {$45$};
  % Bottom bar: split into parts
  \draw[thick,fill=funBlue!20] (0,0) rectangle (3,1);
  \node at (1.5,0.5) {$x$};
  \draw[thick,fill=funBlue!20] (3,0) rectangle (6,1);
  \node at (4.5,0.5) {$x$};
  \draw[thick,fill=funBlue!20] (6,0) rectangle (9,1);
  \node at (7.5,0.5) {$x$};
  \draw[thick,fill=funGreen!20] (9,0) rectangle (10,1);
  \node at (9.5,0.5) {$6$};
  % Brace
  \draw[decorate,decoration={brace,amplitude=5pt,mirror}] (0,-0.2) -- (10,-0.2);
\end{tikzpicture}
\end{center}
\end{questionText}
\correctAnswer{$x = 13$}
\explanation{The model shows $3x + 6 = 45$. Subtract $6$: $3x = 39$. Divide by $3$: $x = 13$.}
\end{practiceQuestion}

\begin{practiceQuestion}{5-2-q30}{gsa}
\begin{questionText}
The flowchart shows the steps used to get from $x$ to $22$. Work backward to find $x$.

\begin{center}
\begin{tikzpicture}[node distance=2cm, auto, scale=0.8]
  \node[draw,rounded corners,fill=funBlue!15,minimum width=1.5cm] (start) {$x$};
  \node[draw,fill=funOrange!15,minimum width=1.5cm,right of=start] (step1) {$\times 5$};
  \node[draw,fill=funOrange!15,minimum width=1.5cm,right of=step1] (step2) {$- 3$};
  \node[draw,rounded corners,fill=funGreen!15,minimum width=1.5cm,right of=step2] (end) {$22$};
  \draw[->,thick] (start) -- (step1);
  \draw[->,thick] (step1) -- (step2);
  \draw[->,thick] (step2) -- (end);
\end{tikzpicture}
\end{center}
\end{questionText}
\correctAnswer{$x = 5$}
\explanation{The equation is $5x - 3 = 22$. Add $3$: $5x = 25$. Divide by $5$: $x = 5$. Check: $5 \times 5 = 25$, $25 - 3 = 22$ \checkmark.}
\end{practiceQuestion}
